\subsection{Single Failure: Recovery Improvements Do Not Ensure Task Gains}
\label{sec:primary-results}
\begin{figure*}[tp]
\centering
\includegraphics[width=\textwidth]{figures/failure-effects.pdf}
\caption{Primary single-failure results.
(a) Absolute mean completion time for all workflows, the selected focus,
the descriptive union of impacted workflows, and the non-focus complement;
small marks show all five seeds.
(b) Descriptive completion ECDFs for 800 task-runs and five focus task-runs
per policy. The 800 observations repeat the same 160 tasks across seeds.
(c,d) Paired failure penalties after subtracting each policy's own F0 control,
for the focus workflow and the whole-workload mean. Negative differences
favor AegisKV. Whiskers use five seed estimates and $df=4$;
neither ECDF supplies additional inferential replicates.}
\Description{Absolute completion times, finite-population distributions,
and paired differences in failure penalty for all five seeds.}
\label{fig:failure-effects}
\end{figure*}


Figure~\ref{fig:workload}(d) records the five matched logical/physical
interventions and their differing injection times and collateral victims.
The descriptive union contains 3--6 impacted workflows per seed. All selected
requests occupy the 64--96K bucket, which contains only 36 of 1,304 calls.
These long-context cases are not representative random failures.

In the failure runs, mean WCT averaged over seeds is 33.836\,s for
Store-only and 33.788\,s for AegisKV. The paired difference is
$-0.048$\,s, with a 95\% interval of $[-0.240,0.143]$\,s.
The mean of the five within-seed P95 differences is $-0.095$\,s
($[-4.831,4.641]$\,s); the corresponding maximum difference is
$-1.734$\,s ($[-16.124,12.655]$\,s). None resolves a consistent
population-wide improvement under the declared admission randomization.
Figure~\ref{fig:failure-effects}(b) retains all 800 task-runs per policy:
328 paired tasks finish sooner with AegisKV, 156 tie, and 316 finish later.
The largest adverse difference is 46.340\,s for workflow
\texttt{db12f9be559698b0} in seed 0; the largest improvement is 57.650\,s.

The selected focus workflows have mean failure-run WCT of 30.502\,s
with Store-only and 34.571\,s with AegisKV. After subtracting each
policy's own failure-free control, the mean focus failure-penalty
difference is $+8.130$\,s ($[-5.116,21.375]$\,s).
Four seed differences are positive and one is negative.
For the whole-workload mean, the same contrast is $+0.088$\,s
($[-0.218,0.395]$\,s). Thus, the small negative absolute WCT difference
does not establish a reduction in the added failure penalty.

\paragraph{Complete-batch completion.}
The complete batch includes admission-slot delay. AegisKV's F1 makespan is
on average 1.022\,s longer ($[-0.871,2.915]$\,s), or a mean paired
relative change of $+0.297$\% ($[-0.252,0.847]$\%).
The makespan failure-penalty contrast is $+0.994$\,s, with interval
$[-1.468,3.456]$\,s. These whole-batch observations likewise do not
establish a consistent completion benefit.

\subsection{Direct Mechanism: Equal Prefixes, Different Recovery Sources}
\label{sec:direct-results}
\begin{figure*}[tp]
\centering
\includegraphics[width=\textwidth]{figures/recovery-mechanism.pdf}
\caption{Direct recovery mechanism for the same primary F1 cases.
(a) Disjoint selected prompt intervals from normal HBM, local Peer-HBM,
Peer-HBM copied over D2D, and Store-DRAM, plus prompt recomputation.
Every bar sums to the frozen focus prompt length. Available tier copies
can overlap; the selected intervals shown here do not.
(b) Recomputed prompt fraction versus recovery-ready latency.
(c) Recovery-ready versus failure-to-request-completion latency.
Lines connect the two policies within seed; numeric labels identify seeds.
Each latency is measured from that execution's actual injection time.
These intervals may overlap and are not additive latency components.}
\Description{Selected recovery source fractions and two paired scatter plots
connecting recomputation, recovery readiness, and request completion.}
\label{fig:recovery-mechanism}
\end{figure*}


All five pairs recover the same effective prompt prefix under both policies:
32,832 tokens in seeds 0, 1, and 4, and 41,024 in seeds 2 and 3.
Consequently, the remaining prompt recomputation is identical within every
pair, accounting for 42.2--51.9\% of the prompt.
Figure~\ref{fig:recovery-mechanism}(a) shows source substitution rather
than additional recovered coverage. Store-only selects Store-DRAM for the
entire recovered prefix. AegisKV uses Peer D2D for seed 0, local Peer-HBM
plus Store for seed 1, almost entirely Store for seeds 2 and 3, and local
Peer-HBM for seed 4. The Peer portion in seeds 2 and 3 is only 64 tokens;
it is retained in the accounting even though it is barely visible at this scale.

Recovery readiness and request completion can move in opposite directions:
seed 2 becomes ready sooner but completes 0.048\,s later after injection;
seed 3 becomes ready later but completes 1.108\,s sooner.
Nor does faster request recovery establish a faster workflow. In seed 1,
the failed request completes 1.509\,s sooner after injection, while the
focus workflow's failure-penalty difference is $+15.985$\,s.
These endpoints have different origins and include different dependencies;
shared scheduling and downstream work mediate their relationship.

Seed 4 further shows why materialization alone is insufficient. Local Peer
adoption materializes the prefix in 1\,ms, compared with 209\,ms for
Store-only. Nevertheless, recovery readiness is later by 0.301\,s and
the focus workflow's failure-penalty difference is $+21.163$\,s.
These observations do not isolate the latency-aware selector's effect from
placement, routing and batching changes.
